Das zentrale Thema des Praxisprojekts ist die Mitwirkung an dem vom Bundesministerium für Bildung und Forschung (BMBF) geförderten Verbundprojekt „GameUp“. Das Ziel dieses Forschungsvorhabens ist die Entwicklung einer innovativen Lösung zur spielerischen Vermittlung von IT-Sicherheitskompetenzen. In diesem Rahmen erfolgt die Arbeit bei der snoopmedia GmbH direkt an der Schnittstelle zwischen technischer Umsetzung und didaktischer Konzeption.

\subsection{Das Projekt GameUp}
Das Projekt „GameUp“ widmet sich der Problemstellung, dass Cyberkriminalität stetig zunimmt, während die Prävention im privaten Umfeld oft an schwer verständlichem Wissen oder mangelnder Motivation der Nutzer scheitert. Das Vorhaben entwickelt hierfür ein Echtzeit-Assistenzsystem in Form eines KI-gestützten Chatbots, der private Verbraucher im digitalen Alltag begleitet. Ein wesentliches Merkmal des Systems ist die proaktive Unterstützung: Nutzer werden gezielt per Push-Benachrichtigung über aktuelle Meldungen im Bereich der IT-Sicherheit informiert. Dies umfasst beispielsweise Warnungen vor neu entdeckten Sicherheitslücken oder akuten Phishing-Wellen.

Die snoopmedia GmbH fungiert dabei als Konsortialführer und ist für die technische Gesamtentwicklung verantwortlich. Weitere Partner wie die TU München und die Universität Paderborn unterstützen das Projekt durch die Bereitstellung didaktischer Inhalte und die Durchführung von Akzeptanzforschungen. Die offizielle Projekt-Website (\url{https://gameup-projekt.de/}) dient dabei als zentrale Informationsplattform für den aktuellen Projektfortschritt.

\subsection{Das Gamification-Konzept}
Ein wesentlicher Pfeiler von „GameUp“ ist der Gamification-Ansatz, der darauf abzielt, die oft als lästig empfundene Beschäftigung mit IT-Sicherheit in ein motivierendes Erlebnis zu transformieren. Das Konzept basiert auf der Metapher einer interaktiven Reise durch eine fiktive Welt, in der verschiedene Orte stellvertretend für spezifische IT-Themen stehen. Die Kernmechanik folgt dabei dem Prinzip „Lernen durch Handeln“:

\begin{itemize}
    \item \textbf{Narrative Einbettung:} Nutzer werden von einem digitalen Guide durch Comic-Szenen geführt, die abstrakte Risiken visualisieren.
    \item \textbf{Interaktive Challenges:} Um neue Orte freizuschalten, müssen Nutzer Aufgaben lösen, wie etwa die Auswahl eines starken Passworts an einem Terminal.
    \item \textbf{Situative Unterstützung:} Durch die Anbindung an aktuelle Sicherheitsmeldungen kann das System den Nutzer genau dann per Push-Nachricht abholen, wenn ein relevantes Ereignis eintritt, und dieses direkt in eine kurze Lerneinheit (Learning Nugget) überführen.
    \item \textbf{Learning Nuggets:} Informationen werden in Form von kurzen, multimedialen Lerneinheiten präsentiert, die der Guide den Nutzern situativ als „Hinweiszettel“ überreicht.
    \item \textbf{Feedback \& Belohnung:} Erfolgreiches Verhalten wird unmittelbar durch den Fortschritt in der Spielwelt sowie durch digitale Abzeichen (Badges) belohnt.
\end{itemize}

Mein persönlicher Aufgabenschwerpunkt im Rahmen des Praxisprojekts liegt in der \textbf{alleinigen technischen Umsetzung} dieses Gamification-Konzepts. Dies umfasst die vollständige softwareseitige Implementierung der spielerischen Mechanismen in die App-Umgebung sowie die technische Integration der Benachrichtigungslogik für die Push-Meldungen, um eine nahtlose Interaktion mit dem Chatbot zu gewährleisten.